% !PRINTED 362
\weberpage{362}
%<<CONT
so dass $x_1, x_2 \ldots x_n$ positive unechte Br\"uche, $a_1, a_2 \ldots a_{n-1}$ ganze positive Zahlen sind.

Nach der Bezeichnung des vorigen Paragraphen ist also $x$ gleich dem Kettenbruch
\[
(a_0, a_1, a_2 \ldots a_{n-1}, x_n).
\]

Diese Kettenbruchentwickelung l\"asst sich aber, wenn $x$ irrational ist, \gesperrt{unbegrenzt fortsetzen}, d. h. wir k\"onnen $n$ beliebig gross annehmen. Es entsteht der folgende aus dem vorangehenden dadurch, dass man
\[
x_n = a_n + \frac{1}{x_{n+1}}
\]
setzt. Auch hier nennen wir die ganzen Zahlen $a_0, a_1, a_2 \ldots a_{n-1}$ die \gesperrt{Theilnenner}, $x_n$ die \gesperrt{Schlusszahl}.

Die Kettenbruchentwickelung der rationalen Zahlen ist als Specialfall darin enthalten. Der Kettenbruch kann (nach Satz I.) h\"ochstens um ein Glied fortgesetzt werden, wenn die Schlusszahl eine ganze Zahl ist.

\begin{center}
\S.\ 117.

Die N\"aherungsbr\"uche.
\end{center}

Wenn wir einen Kettenbruch
\weq{1}{x = (a_0, a_1, a_2 \ldots a_{n-1}, x_n),}
worin wir $x_n$ als eine variable Gr\"osse ansehen k\"onnen, in einen gew\"ohnlichen Bruch verwandeln, so erhalten wir im Z\"ahler und im Nenner einen linearen Ausdruck in $x_n$. Es wird also
\weq{2}{x = \frac{P_n x_n + P_{n-1}}{Q_n x_n + Q_{n-1}},}
worin $P_n$, $Q_n$, $P_{n-1}$, $Q_{n-1}$ ganze rationale Functionen der $a_0, a_1, a_2 \ldots a_{n-1}$ sind. Dass dies richtig ist, ergiebt sich f\"ur die ersten Werthe $n = 1, 2 \ldots$ durch unmittelbare Rechnung:
\[
\begin{aligned}
x &= a_0 + \frac{1}{x_1} = \frac{a_0 x_1 + 1}{x_1}\\
&= a_0 + \cfrac{1}{a_1 + \cfrac{1}{x_2}} = \frac{(a_0 a_1 + 1)x_2 + a_0}{a_1 x_2 + 1}.
\end{aligned}
\]
